%% kapa: exam=1 round=29 year=2018 period=1 form=A subject=real_estate subject_ko=부동산학원론 seq=10 number=90
%% source: Q-Net 감정평가사 기출문제 (cst00309) artlSeq=5202808 file=extracted/1008/5202808/2018년 제29회 감정평가사 1차시험 기출문제/2018년 제29회 감정평가사 1차시험 기출문제/1교시 A형.pdf
%% conversion: kordoc 4.24.51 -> tools/gen1.py (auto); review=proofread-sonnet; flags=none
\begin{question}{90}
부동산의 가치와 가격에 관한 설명으로 옳지 않은 것은?
\begin{choices}
  \item 일정시점에서 부동산가격은 하나 밖에 없지만, 부동산가치는 여러 개 있을 수 있다.
  \item 부동산가격은 장기적 고려 하에서 형성된다.
  \item 부동산의 가격과 가치 간에는 오차가 있을 수 있으며, 이는 감정평가 필요성의 근거가 된다.
  \item 부동산가격은 시장경제에서 자원배분의 기능을 수행한다.
  \item 부동산가치는 부동산의 소유에서 비롯되는 현재의 편익을 미래가치로 환원한 값이다.
\end{choices}
\end{question}
\answer{5}
